\documentclass{article}

\usepackage{amsmath, amssymb, amsthm}
\usepackage[a4paper, margin=1in]{geometry}

\title{PRIMES in P : The Question}
\author{sid}
\date{\today}

\newtheorem{theorem}{Theorem}
\newtheorem{lemma}{Lemma}

\newcommand{\Z}{\mathbb{Z}}

\begin{document}

\maketitle

\section{The Problem of PRIMES}

Let $n \in \mathbb{Z}$ with $n > 1$. Determine whether $n$ is \textbf{prime}.

\section{Input Model}

For an integer $n$.
\[
\text{Length of n in binary} = k = \lfloor \log_2 n \rfloor + 1
\]

$\implies$ the complexity of algorithms will be measured as a function of $k$, not $n$.

\section{Trial Division is a Lie}

The naivest approach would be to just divide $n$ by every integer from $1$ to $n$.
This should take $O(\sqrt{n})$ time.

However,
$$
O(\sqrt{n})=O(n^{1/2})=O({2^{(k-1)}}^{1/2})=O(2^{\sqrt{k}})
$$

This is \textbf{EXP} time.

Thus, the key question becomes:

\begin{center}
\textit{Is there a deterministic algorithm that decides primality in time polynomial in $k$?}
\end{center}
i.e.,
\begin{center}
\textit{Is PRIMES in P?}
\end{center}


\end{document}